% Generated draft: requires mathematical and source review.
\documentclass[11pt]{article}
\usepackage[margin=25mm]{geometry}
\usepackage{fontspec}
\setmainfont{Latin Modern Roman}
\usepackage{amsmath,amssymb,mathtools,unicode-math}
\setmathfont{Latin Modern Math}
\usepackage{enumitem,needspace,xurl,hyperref}
\setlength{\parindent}{0pt}
\setlength{\parskip}{5pt}
\setlength{\emergencystretch}{3em}
\newcommand{\N}{\mathbb{N}}
\newcommand{\Z}{\mathbb{Z}}
\newcommand{\Q}{\mathbb{Q}}
\newcommand{\R}{\mathbb{R}}
\newcommand{\C}{\mathbb{C}}
\newcommand{\F}{\mathbb{F}}
\newcommand{\A}{\mathbb{A}}
\newcommand{\PP}{\mathbb{P}}
\newcommand{\E}{\mathbb{E}}
\newcommand{\eps}{\varepsilon}
\newcommand{\dd}{\,\mathrm{d}}
\begin{document}
\section*{UnsolvedMath ID: 30006011}
\textbf{Source problem code:} catalogue entry
\subsection{Definitions and mathematical statement}
\textbf{Scope and missing definitions.} This record is a contextual continuation, not a recoverable standalone conjecture. The supplied registry marks its statement as unrecoverable and identifies the conjecture recorded under source code \texttt{OWR-14298588-010} as its intended context. Properties (1), (2), and (3), the preceding construction, and the relevant definitions are absent.

Retain the source symbols \(G\), \(E\), \(\ell\), \(\iota\), and \(\psi\), without supplying unstated hypotheses. The excerpt calls \(\psi\) an A-parameter with respect to \(\iota\), but does not define that notion or the domains of these data. It calls \(\psi\) tempered when its restriction to the second \(\mathrm{SL}_2\) is trivial; the parameter's full domain is not supplied. The notation \(D_{\mathrm{lis}}(\operatorname{Bun}_G,\Q_\ell)\) below follows the transcription; the underlying geometric setting, category, and precise coefficient conventions require the missing context.

\textbf{Conditional assertion retained from the excerpt.} The source states that the categorical local Langlands conjecture would imply the existence of an object
\[
0\ne\mathcal F_\psi\in D_{\mathrm{lis}}(\operatorname{Bun}_G,\Q_\ell)
\]
satisfying its property (1). Since property (1), the precise categorical conjecture, and the admissible parameters are not supplied, this is an incomplete conditional assertion, not a fully specified existence problem.

\textbf{Spectral evidence and qualifications.} The excerpt separately labels as a theorem an assertion that, for an A-parameter \(\psi\) with respect to \(\iota\), there exists a nonzero ind-coherent sheaf satisfying an eigenrelation. Its displayed spectral space and eigenrelation are corrupted in extraction. The surviving notation involves \(Z^1\), \(W_E\), a dual-group symbol, a quotient by that group, coefficient notation involving \(\Q_\ell\), and an isomorphism involving \(V\), components \(V_i\), tensor products with the sheaf, a sum indexed by \(i\), and shifts \([-i]\). The space, summation, representations, grading, tensor action, and quantification over \(V\) cannot be specified reliably from this record. Consequently, no reconstructed formula is asserted here as the theorem's precise statement.

For tempered \(\psi\), the source describes the preceding ind-coherent construction as essentially obtained from the regular representation of the centralizer of \(\psi\), referring to [8, Section 3.1]. In general it attributes inspiration to [1], which it describes as comparing categorical local Langlands with \(p\)-adic Adams--Barbasch--Vogan theory in the unipotent case through Koszul duality. The authors present this as motivation for property (3), and describe purely spectral formulations and verification of properties (2) and (3) as work in progress. They report a check for \(\mathrm{GL}_2\), without supplying enough context to identify its exact assertion. They explicitly require some version of classical local Langlands before property (3) makes sense, because Adams--Barbasch--Vogan packets require that input.

The source also relates pushforward along the Hodge--Tate period map from a perfectoid compact Shimura variety, more precisely from the Igusa stack, to eigensheaves. It associates the modified eigenproperty (1) with [6, 7] and [2], describing its interpretation in [2] as a special case of shearing. None of these comments supplies the missing properties or a separate precise conjecture.
\subsection{Short English statement}
This excerpt discusses consequences and spectral evidence for an Arthur-parameter eigensheaf conjecture recorded elsewhere. The defining properties and hypotheses are missing, and the displayed spectral theorem is corrupted. It therefore supports only an explicitly incomplete contextual statement requiring human mathematical review. The supplied open label is a registry assessment, not an independently verified mathematical conclusion.
\subsection{Sources}
\begin{itemize}
\item Supplied local UnsolvedMath registry record: permanent ID \(30006011\), exact source code \texttt{catalogue entry}, title ``Arthur-Eigensheaf Spectral-Side Continuation.'' The registry gives proposed year 2024, no named proposer, status ``open,'' and statement status ``unrecoverable.'' Its supplied verification note identifies this passage as discussion and evidence following the conjecture recorded under \texttt{OWR-14298588-010}, rather than a separate conjecture; this verification note has not been independently checked here.
\item Supplied report citation: \emph{Arithmetic Geometry (2025)}, Oberwolfach Reports, pp. 1855--1912, DOI \url{https://doi.org/10.4171/owr/2024/33}. The report citation and link provide provenance, but the supplied excerpt does not identify its contribution author, precise passage location, or missing definitions.
\item Supplied literature source ``Hansen--Mann'': \url{https://arxiv.org/abs/2606.00983}. No title, specific result, or identification with the excerpt's numbered references is supplied.
\item Literature assessment dated exactly \texttt{2026-08-22}: supplied dated curation, not independently verified. It labels the record open, characterizes it as contextual commentary rather than a standalone proved claim, and states that the general categorical realization invoked remains unresolved. The supplied background requests independent verification of time-sensitive status and cited claims.
\item Internal references [1], [2], [6], [7], and [8, Section 3.1] occur in the excerpt, but their bibliographic entries are absent. No complete authoritative citation establishing a standalone formulation is supplied.
\item Attribution retained from the excerpt: ``TK'' was supported by JSPS KAKENHI Grant Number \texttt{24K16895}. The supplied record does not expand these initials or identify the speakers of its first-person statements.
\end{itemize}
\end{document}
